% Generated draft: requires mathematical and source review.
\documentclass[11pt]{article}
\usepackage[margin=25mm]{geometry}
\usepackage{fontspec}
\setmainfont{Latin Modern Roman}
\usepackage{amsmath,amssymb,mathtools,unicode-math}
\setmathfont{Latin Modern Math}
\usepackage{enumitem,needspace,xurl,hyperref}
\setlength{\parindent}{0pt}
\setlength{\parskip}{5pt}
\setlength{\emergencystretch}{3em}
\newcommand{\N}{\mathbb{N}}
\newcommand{\Z}{\mathbb{Z}}
\newcommand{\Q}{\mathbb{Q}}
\newcommand{\R}{\mathbb{R}}
\newcommand{\C}{\mathbb{C}}
\newcommand{\F}{\mathbb{F}}
\newcommand{\A}{\mathbb{A}}
\newcommand{\PP}{\mathbb{P}}
\newcommand{\E}{\mathbb{E}}
\newcommand{\eps}{\varepsilon}
\newcommand{\dd}{\,\mathrm{d}}
\begin{document}
\section*{UnsolvedMath ID: 30002601}
\textbf{Source problem code:} catalogue entry
\subsection{Definitions and mathematical statement}
The supplied record asks whether two coprime one-variable exponential sums have finitely many common zeros. It does not specify the coefficient field, the domain of the variable, or the ring in which divisibility is taken. The following is an explicitly provisional complex-analytic formulation requiring mathematical review.

Let
\[
\mathcal E=\left\{z\longmapsto\sum_{j=1}^{m}a_j e^{\lambda_j z}:m\in\Z_{\geq 0},\ a_j,\lambda_j\in\C\right\},
\]
where the empty sum is zero and the functions are defined on \(\C\). Regard \(\mathcal E\) as a ring under pointwise operations. A function \(d\in\mathcal E\) divides \(f\in\mathcal E\) if \(f=dh\) for some \(h\in\mathcal E\). A unit is a function \(d\in\mathcal E\) for which \(dv=1\) for some \(v\in\mathcal E\).

Under these provisional conventions, the proposed statement is
\[
\forall f,g\in\mathcal E,\qquad
\left[\text{every common divisor of }f\text{ and }g\text{ in }\mathcal E\text{ is a unit}\right]
\Longrightarrow
\#\{z\in\C:f(z)=g(z)=0\}<\infty.
\]
The supplied clean statement says ``no nonconstant common divisor,'' whereas the original statement says ``no nontrivial common divisor.'' These expressions require reconciliation: in \(\mathcal E\), every \(e^{\lambda z}\) is a unit, including nonconstant ones, and divides every element. Thus the literal exclusion of nonconstant common divisors cannot express coprimality in this ring. The formulation above interprets ``nontrivial'' as ``nonunit''; that interpretation and the definition of exponential sum remain unverified.

The original record also mentions finiteness of \((A\setminus A')\cap G\) under unspecified finite-rank and anomalous-subvariety hypotheses. Its supplied statement-verification note assigns that separate higher-dimensional Mordell--Lang generalization to \texttt{OWR-12985-004}; it is therefore not included in the proposed statement for this record.
\subsection{Short English statement}
Under an explicitly provisional interpretation using finite sums of complex exponentials on \(\C\), the conjecture asks whether two functions with no common nonunit exponential-sum divisor have only finitely many common zeros. The domain and divisor convention need review. The supplied registry labels the problem open; that status has not been independently verified.
\subsection{Sources}
\begin{itemize}
\item Local UnsolvedMath registry, supplied complete record for permanent ID \(30002601\), source code \texttt{catalogue entry}, title \emph{Finiteness of Common Zeros of Coprime Exponential Sums}, collection \emph{Oberwolfach Reports Open Problems}. The proposed year is 2014; no individual proposer is supplied.
\item Supplied report citation: \emph{Okounkov Bodies and Applications (2015)}, Oberwolfach Reports, pp.~1459--1513, DOI \texttt{10.4171/owr/2014/27}. \url{https://doi.org/10.4171/owr/2014/27}. The citation's 2015 date differs from the report's supplied 2014 designation; this discrepancy is unresolved.
\item Supplied literature source: \emph{Oberwolfach Report 27/2014}. \url{https://ems.press/content/serial-article-files/46517}.
\item Supplied literature source: Heittokangas et al., exponential polynomials survey. \url{https://londmathsoc.onlinelibrary.wiley.com/doi/full/10.1112/blms.12719}.
\item Supplied literature source: Kazarnovskii, \emph{On exponential algebraic geometry}. \url{https://www.mathnet.ru/php/getFT.phtml?jrnid=rm&option_lang=eng&paperid=10184&what=fullteng}.
\item Literature assessment dated exactly \texttt{2026-08-21}: the supplied curation labels the conjecture open, describes it as far from proved in current primary surveys, and states that special cases follow from Mordell--Lang or Skolem--Mahler--Lech type results. These are supplied dated curation claims, not independently verified literature findings.
\item The supplied registry labels the statement \texttt{corrected\_verified} and reports that the exponential-sum conjecture is separate from the higher-dimensional generalization assigned to \texttt{OWR-12985-004}. This source-verification claim has not been independently checked.
\end{itemize}
\end{document}
